\documentclass[12pt, a4paper]{academics}

\academicsCourse{计算机系统基础}{Introduction to Computer Systems}
\academicsReport{实验课程总结}{2026 年 5 月 25 日}
\academicsArthur{华博文}{19240212}

\begin{document}

\makeacademicscover

\section{实验内容}

本学期共有七个实验，内容从基础环境安装逐步推进到数据表示、编译调试和逆向分析，整体难度是递进的。

\subsection{实验一：环境安装与测试}

实验一主要熟悉课程所需工具链，验证 \mintinline{text}{gcc}、\mintinline{text}{objdump} 和基础编译流程是否可用。通过编译简单的 \mintinline{text}{hello.c}，我第一次完整接触了预处理、编译、汇编、链接到可执行文件的全过程，也认识到源代码和目标文件之间存在明显差异。

\subsection{实验二：位运算与补码}

实验二围绕位运算和补码展开，重点是理解按位与、或、异或、移位以及有符号整数的机器级表示。这个实验让我更加熟悉二进制数在底层是如何参与逻辑运算和算术运算的，也让我意识到类型转换和位宽差异会直接影响结果。

\subsection{实验三：整数数据的表示}

实验三进一步研究整数在不同位宽下的表示方法，包括无符号数、带符号数和补码表示。通过分析十六进制数值与真值之间的对应关系，我对整数溢出、符号扩展和机器数解释方式有了更清楚的认识。

\subsection{实验四：浮点数的表示}

实验四重点分析 IEEE-754 单精度浮点数的结构，包括符号位、阶码和尾数，以及特殊值如 \mintinline{text}{inf}、\mintinline{text}{nan}、\mintinline{text}{+0}、\mintinline{text}{-0} 和非规格化数。这个实验最直接的体会是：浮点数并不是十进制小数的精确映射，很多看似相近的数在机器里会落到相同或相邻的表示上。

\subsection{实验五：程序的编辑编译调试与运行}

实验五让我把多个源文件组织成一个完整程序，并使用 \mintinline{text}{clang}、\mintinline{text}{objdump} 和 \mintinline{text}{gdb} 进行编译、反汇编和调试。这个实验最重要的收获是学会了从工程角度理解程序：函数如何分文件组织，链接如何把各部分拼接起来，以及如何通过调试器观察程序执行过程。

\subsection{实验六：数据的存储}	

实验六围绕内存排列方式、大小端存储和多类型数据解释展开。通过分析同一组字节如何被解释成无符号整数、补码、浮点数、BCD 码、ASCII 字符和汉字编码，我对“数据在内存中只是字节序列”这件事有了更深的理解，也更能从底层视角看待数据结构和编码。

\subsection{实验七：二进制炸弹拆弹实验}

实验七是本学期最综合的一次实验。它要求我把前面学到的编译、反汇编、数据表示、栈帧、递归、链表和条件跳转等知识综合起来，借助 \mintinline{text}{objdump} 和 \mintinline{text}{gdb} 逐步推导每一阶段的正确输入。这个实验让我真正体会到逆向分析不是“猜答案”，而是沿着程序执行路径一点点恢复出程序的意图。

\section{心得体会}

这七个实验从“会用工具”逐步过渡到“理解程序为什么这样运行”。前几个实验帮助我夯实了位运算、整数表示、浮点表示和内存存储的基础；后几个实验则把这些知识落实到编译、调试和逆向分析上。相比单纯背概念，这种边做边分析的方式让我更容易把分散的知识连成一条线。

我最大的体会是，计算机系统基础不是孤立的知识点集合，而是一套层层衔接的体系：数据如何表示，程序如何编译，指令如何执行，内存如何组织，调试工具如何观察这些过程。以后遇到更复杂的问题，我会先从数据表示和执行路径入手，再结合工具去验证推理，而不是只停留在表面现象。

\section{结论}

本学期的七个实验把我从基础工具使用带到了系统级分析的层面，也让我逐步建立起对“程序如何被机器执行”的整体认识。最初我更多是在完成命令和现象观察，后来逐渐能够主动分析数据表示、代码结构和执行流程，这种变化比单次实验结果更重要。

同时，我也发现自己在分析复杂程序时还不够熟练，尤其是在把汇编逻辑和高级语言结构互相对应起来时，需要更多训练。后续如果继续接触类似实验，我希望自己能更快定位关键函数、快速整理推理链条，并把分析过程写得更清楚、更简洁。

\end{document}
